% ============================================================
% preamble.tex — book-production ENGINE (brand-agnostic)
% ============================================================
% Everything in this file is structural: page layout, table
% mechanics, front/main/back-matter page-numbering control,
% landscape-page handling, header/footer behavior. None of it is
% specific to any one brand or course — it's meant to
% be reused unmodified across every book built on this toolkit.
%
% Colors, fonts, and the book title are NOT defined here — they live
% in brand.tex, loaded before this file (see _quarto.yml's
% `include-in-header:` order). This file references color/font names
% (giPrimary, \headingfont, \giBookTitle, etc.) that brand.tex is
% responsible for defining. If you're adapting this toolkit for a
% different brand, brand.tex is what you replace — this file
% shouldn't need to change.
% ============================================================

% ---- Layout essentials (images, landscape pages) ----
\usepackage{graphicx}
\usepackage{pdflscape}

% ============================================================
% HEADING STYLES (the h1/h2/h3 rules of the LaTeX world)
% ============================================================
\usepackage{titlesec}
% \raggedright on the title text (both here and in the 4th
% argument below) prevents a real, confirmed typesetting defect: a
% long chapter title that wraps to two lines otherwise gets fully
% JUSTIFIED at \Huge size — LaTeX's default paragraph behavior — and
% stretching a handful of huge words to exactly fill a line's width
% produces visibly uneven, overly wide gaps between words on the
% wrapped line. Reproduced directly with a real long title
% ("Dissecting Twelve Unique Security Failures and Governance Layers
% Across a Decade") and confirmed \raggedright fixes it with zero
% effect on short, single-line titles. Parts (below) don't need this
% same fix — \centering, which they already use, doesn't stretch
% inter-word spacing the way justification does.
\titleformat{\chapter}[display]
  {\headingfont\Huge\bfseries\color{giPrimary}\raggedright}
  {\chaptertitlename\ \thechapter}{20pt}{\Huge\raggedright}
  [\color{giCrimson}\titlerule]
% Parts sit one level above chapters — their own divider page, larger
% and centered, distinct from a chapter opening rather than just a
% bigger version of one. \thepart defaults to roman numerals (Part
% I, Part II); switch to \arabic{part} in this block if you'd rather
% have "Part 1, Part 2." Only relevant if you use `part:` entries in
% _quarto.yml's chapters: list — harmless and inert if you don't; see
% README's "Parts" section for the full syntax and behavior.
%
% \titleclass{\part}{top} switches \part from titlesec's default
% "page" class to "top" (the same class \chapter uses) BEFORE
% \titleformat/\titlespacing configure its appearance below. This
% matters because titlesec fully replaces LaTeX's own \part command
% with its own \ttl@pageclass mechanism — the kernel's \@endpart
% (which normally forces a fresh page after every part title, even
% with nothing else on the page) is never even called once titlesec
% is loaded, so redefining \@endpart directly has no effect; titlesec
% controls this via its own page classification instead, and "page"
% class is what forces that automatic post-title page break. "top"
% still starts the part on its own fresh page (same as \chapter
% does) but doesn't force whatever follows onto a SEPARATE page —
% needed so a short part-overview paragraph placed right after a
% `# Part Title` heading lands on the same page as the title instead
% of being stranded alone on an otherwise-blank page.
\titleclass{\part}{top}
\titleformat{\part}[display]
  {\centering\headingfont\Large\bfseries\color{giCrimson}}
  {\partname\ \thepart}{16pt}{\huge\color{giPrimary}}
  [\color{giCrimson}\titlerule]
\titlespacing*{\part}{0pt}{80pt}{60pt}
\titleformat{\section}
  {\headingfont\Large\bfseries\color{giPrimary}}{\thesection}{1em}{}
\titleformat{\subsection}
  {\headingfont\large\bfseries\color{giSecondary}}{\thesubsection}{1em}{}
\titlespacing*{\chapter}{0pt}{0pt}{20pt}
\titlespacing*{\section}{0pt}{18pt}{6pt}

% ---- Body text color ----
% \AtBeginDocument defers this until the full preamble has loaded —
% needed specifically for the Quarto-extension packaging (see
% "Installing via Quarto extension" in README), where this file
% loads BEFORE brand.tex (the extension mechanism forces that order,
% the opposite of what a plain include-in-header list would do) — so
% \color{giInk} would otherwise fire before giInk is even defined.
% Confirmed via an actual render, not assumed: this was the ONLY line
% in this file that executes a brand color directly rather than
% storing it inside a macro/titleformat definition for later use, so
% no other line needed the same fix.
\AtBeginDocument{\color{giInk}}

% ---- Body text spacing ----
\usepackage{setspace}
\setstretch{1.15}   % line-height equivalent

% ============================================================
% FRONT/MAIN MATTER CONTROL
% Pandoc's book template unconditionally inserts its own \mainmatter
% immediately before the "chapters:" content begins (i.e. right
% before Copyright), regardless of any \frontmatter the user places
% around it. That silently forces arabic numbering on Copyright/
% Preface/TOC. We neutralize that one automatic call and expose our
% own \giEnableMainMatter switch, fired explicitly once — at the end
% of contents.qmd, right before the Introduction chapter — which is
% the only place \mainmatter should actually take effect.
% ============================================================
\makeatletter
\newif\ifgi@realmainmatter
\gi@realmainmatterfalse
\let\gi@origmainmatter\mainmatter
\renewcommand{\mainmatter}{%
  \ifgi@realmainmatter\gi@origmainmatter\fi
}
\newcommand{\giEnableMainMatter}{\gi@realmainmattertrue\mainmatter\gi@inchapterstrue}

% Records the roman page count reached at the end of front matter
% (the TOC page), so back matter can resume roman numbering from
% there — e.g. front matter ends at "iii", back matter picks up at
% "iv" — rather than restarting at "i" and producing two different
% pages both labeled "i" in the same document.
\newcounter{giFrontMatterLastPage}
\newcommand{\giMarkFrontMatterEnd}{\setcounter{giFrontMatterLastPage}{\value{page}}}
\newcommand{\giEnableBackMatter}{%
  \backmatter
  \pagenumbering{roman}%
  \setcounter{page}{\numexpr\value{giFrontMatterLastPage}+1\relax}%
  \gi@inchaptersfalse
}

% ============================================================
% APPENDICES
% \giEnableAppendix switches from numbered chapters (Chapter 1, 2,
% 3...) to lettered appendices (Appendix A, B, C...) — LaTeX's
% standard \appendix mechanism, called at the point in the book where
% appendices should start (after the last numbered chapter, before
% the first appendix — same "must fire at the end of the PREVIOUS
% file" placement rule as \giEnableMainMatter/\giEnableBackMatter
% elsewhere in this project, since Quarto moves raw LaTeX found
% before a heading to after it).
%
% This needed surprisingly little custom code — confirmed by actually
% building and testing a 3-appendix example (including one using
% \giLandscapeChapterOpener, since landscape content is common in
% appendices) rather than assuming it would just work:
%   - Chapter titles automatically show "Appendix A" instead of
%     "Chapter 3" — no changes needed to \titleformat{\chapter} or
%     \giLandscapeChapterHeader, since both already reference
%     \chaptertitlename (a standard LaTeX kernel command that tracks
%     appendix mode on its own) rather than hardcoding "Chapter".
%   - Table/figure numbering automatically follows suit ("Table A.1",
%     "Figure B.1"), since it derives from \thechapter, which
%     \appendix redefines to use letters instead of digits.
%   - The List of Tables/Figures "chapters only" gate
%     (\ifgi@inchapters, see below) is UNCHANGED by \giEnableAppendix
%     on purpose — appendix content should still count for those
%     lists, exactly like a regular numbered chapter does, and it
%     already does automatically since nothing about the gate itself
%     is tied to whether \c@chapter is numeric or lettered.
%   - Page numbering is untouched by \appendix itself (stays arabic,
%     continuing the same sequence as the preceding chapters) —
%     \giEnableBackMatter afterward is what switches to roman for the
%     back matter that follows the appendices, exactly as it already
%     did before appendices existed.
%
% Usage — at the end of the last numbered chapter's file:
%   ```{=latex}
%   \giEnableAppendix
%   ```
% Then write each appendix like a normal chapter (or, for a
% landscape-only appendix, using \giLandscapeChapterOpener exactly as
% documented above) — no other syntax changes needed. Move
% \giEnableBackMatter to the end of the LAST appendix file instead of
% the last numbered chapter, since back matter now follows the
% appendices instead of following the chapters directly.
% ============================================================
\newcommand{\giEnableAppendix}{%
  \appendix
}

% ============================================================
% LIST OF FIGURES / LIST OF TABLES — CHAPTERS ONLY, ENFORCED
% Pandoc treats any standalone Markdown image with alt text as a
% captioned figure by default — so a completely normal-looking image
% insertion in a front- or back-matter file (a logo on the "About"
% page, say) silently ends up in the List of Figures too, since nothing
% about \caption{} itself knows or cares what part of the book it's
% in. Relying on "just don't caption images outside chapters" as a
% writing convention is too easy to violate by accident (real
% incident: an ordinary ![...](...) image on a back-matter page did
% exactly this). Fixed here at the mechanism level instead: \giEnableMainMatter
% flips \ifgi@inchapters on, \giEnableBackMatter flips it back off, and
% \addcontentsline is wrapped (further down, AFTER hyperref — see the
% comment there for why the order matters) so lof/lot entries are
% silently dropped whenever that flag is off — regardless of how the
% image or table got captioned. Front/back matter can caption images
% and tables freely now; they just won't appear in the List of
% Figures/Tables. Everything else (\addcontentsline{toc}{...} for
% chapter/section headings, etc.) passes through untouched.
% ============================================================
\newif\ifgi@inchapters
\gi@inchaptersfalse
\makeatother

% ---- Link colors ----
\usepackage{hyperref}
\hypersetup{
  colorlinks=true,
  linkcolor=giAccent,
  urlcolor=giAccent
}

% This must come AFTER \usepackage{hyperref}, deferred via
% \AtBeginDocument AND re-applied via \AddToHook{package/caption/after}
% — belt and suspenders, because Quarto/Pandoc loads the `caption`
% package LAZILY (only at the first table/figure that actually needs
% it, wherever that falls in the document), and that package fully
% takes over \caption's internals — including, it turns out, keeping
% its own internal reference to whatever \addcontentsline looked like
% when IT loaded, rather than deferring to whatever the command
% currently means. That makes patching \addcontentsline itself a dead
% end: it compiles cleanly and silently does nothing, because
% `caption`'s own internal figure/table-caption code path never
% actually calls the patched version. \AddToHook{package/caption/after}
% is the LaTeX kernel's built-in mechanism for "run this code
% immediately after package X finishes loading, no matter when that
% turns out to be" — used here to patch \caption itself instead, once
% `caption` package's own version of it is the thing in effect.
%
% The actual technique: \caption[<short>]{<full>} is the standard,
% universally-supported LaTeX idiom where the OPTIONAL argument
% controls the .lof/.lot entry and the mandatory argument controls
% what's shown on the page — they're independent. Passing an EMPTY
% optional argument (\caption[]{...}) keeps the full caption visible
% on the page while adding nothing to the list. So this rewrites
% \caption{text} into \caption[]{text} whenever \ifgi@inchapters is
% false — outside a chapter's content, figures/tables keep their
% on-page caption (still numbered, since the underlying figure/table
% counter is a single global sequence LaTeX doesn't reset between
% front/main/back matter — a real but minor cosmetic wrinkle, not a
% correctness issue) but are correctly excluded from the List of
% Figures/Tables either way. This project's tables/figures only ever
% use the plain single-argument \caption{text} form (never a manual
% custom short caption via \caption[short]{text}), so that's the only
% form this gate needs to handle — it isn't meant to support the two-argument
% form.
\makeatletter
\newif\ifgi@captionpatched
\gi@captionpatchedfalse
\newcommand{\gi@patchcaption}{%
  \ifgi@captionpatched\else
    \gi@captionpatchedtrue
    \let\gi@origcaption\caption
  \fi
  \renewcommand{\caption}[1]{%
    \ifgi@inchapters
      \gi@origcaption{##1}%
    \else
      \gi@origcaption[]{##1}%
    \fi
  }%
}
\AtBeginDocument{\gi@patchcaption}
\AddToHook{package/caption/after}{\gi@patchcaption}
\makeatother

% ============================================================
% HEADER / FOOTER
% One item per header line — alternating between the book title and
% the current chapter name by page — instead of both simultaneously
% at left+right. Two items sharing a line risks visual collision or
% an awkward wrap when either title is long; alternating keeps each
% header line short and gives the book title equal visibility to the
% chapter name across the book, without doubling up.
%
% Chapter-opening pages intentionally get NO header (LaTeX's default
% \pagestyle{plain} — page number only): the big chapter title
% directly below already states what page you're on, so repeating it
% in a header on the same page is redundant.
% ============================================================
\usepackage{fancyhdr}
\pagestyle{fancy}
\fancyhf{}
% Wrapped in a fixed-width, right-aligned box so a long chapter name
% or title wraps to a second line instead of overflowing the margin
% or colliding with anything else (nothing else occupies this line).
\setlength{\headheight}{27pt}
\fancyhead[R]{%
  \parbox[b]{0.85\textwidth}{\raggedleft\small\headingfont\color{giSecondary}%
  \ifodd\value{page}\leftmark\else\giBookTitle\fi}%
}
\fancyfoot[C]{\small\thepage}
\renewcommand{\headrulewidth}{0.4pt}
\renewcommand{\footrulewidth}{0pt}

% ============================================================
% TABLE STYLING
% ============================================================
\usepackage{booktabs}
\usepackage{array}
\setlength{\arrayrulewidth}{0.6pt}   % border thickness
\renewcommand{\arraystretch}{1.3}    % row height / cell padding

% Pandoc renders every Markdown table (including narrow ones) inside
% a `longtable`, which centers itself on the page by default — not
% the standard convention for body-text tables. Left-align instead,
% flush with the surrounding paragraph text (professional norm; also
% matches how the images/headings in this template are set).
\setlength\LTleft{0pt}
\setlength\LTright{0pt plus 1fil}

% Usage in a chapter file, right before a table:
%   \rowcolors{2}{tableAltRow}{white}   <- zebra-stripes body rows,
%                                           leaves the header row plain
%                                           (distinguished by booktabs
%                                           rules above/below it)
%
% Note: per-row \rowcolor commands only work inside hand-written raw
% LaTeX tables — Pandoc's auto-generated tables (from Markdown pipe
% syntax) don't expose a hook to color just the header row. If you
% need a shaded header specifically, that table needs to be written
% as raw LaTeX rather than a Markdown table — ask if you hit this.
% See chapter-01.qmd's landscape example for a full worked version
% with header/first-column/alternating-row shading combined.

% ============================================================
% LANDSCAPE PAGE HEADER / FOOTER
% pdflscape rotates the *body* content of a `landscape` block so it
% displays right-way-up, but it has no effect on \fancyhead/\fancyfoot
% — those are drawn against the physical (portrait) page edges, so on
% a landscape page they end up rotated sideways along the margin.
% Fix: don't use fancyhdr's per-page hooks on landscape pages at all.
% Suppress them (\thispagestyle{empty}) and instead print a matching
% header/footer line as ordinary content, inside the `landscape`
% block itself — so it's part of the same rotated box as the table
% and reads normally horizontal, top and bottom of the page, exactly
% like every portrait page.
% ============================================================
\newcommand{\giLandscapeHeader}{%
  \thispagestyle{empty}%
  {\small\headingfont\color{giSecondary}\leftmark\hfill\giBookTitle\par}%
  \vspace{2pt}{\color{giCrimson}\hrule height 0.4pt}\vspace{10pt}%
}
\newcommand{\giLandscapeFooter}{%
  \vfill\hfill{\small\thepage}\par
}

% ============================================================
% LANDSCAPE-ONLY CHAPTER OPENER
% A normal chapter (via a Markdown "# Heading") always produces its
% own full portrait page: \clearpage, then the chapter title typeset
% with \titleformat{\chapter}'s styling. That's fine when real
% portrait content follows it on the same page — but if a chapter's
% first (or only) content is a landscape table/figure, the normal
% flow wastes an entire page: the portrait title page ends up nearly
% blank (title only, nothing else fits before the landscape rotation
% has to start its own page), and the actual content only begins on
% the page after that.
%
% \giLandscapeChapterOpener{Title} is a drop-in replacement for the
% Markdown heading, used ONLY for a chapter whose first content is
% landscape — it does the same bookkeeping a normal \chapter does
% (advances the chapter counter, adds the Table of Contents entry,
% sets the running header text) but deliberately WITHOUT typesetting
% a separate portrait title page. Pair it with
% \giLandscapeChapterHeader (no arguments — it reuses the title
% \giLandscapeChapterOpener already stored) INSIDE the landscape
% block instead of the normal \giLandscapeHeader, so the chapter
% title itself appears as part of the landscape page's own header,
% horizontal, at the top — not as a separate page.
%
% Usage (raw LaTeX, not a Markdown heading — see chapter-01.qmd for
% a full working example):
%   ```{=latex}
%   \giLandscapeChapterOpener{A Chapter That Is Only A Landscape Table}
%   \begin{landscape}
%   \giLandscapeChapterHeader
%   ... table or figure ...
%   \giLandscapeFooter
%   \end{landscape}
%   ```
%
% If the chapter has ANY portrait content before the landscape block
% (even one paragraph), use a normal Markdown heading and
% \giLandscapeHeader instead, exactly as documented elsewhere — this
% pair is specifically for the landscape-content-only case.
%
% One more real thing this needed: a chapter's .qmd file with NO
% Markdown heading at all (true for a landscape-only chapter, since
% \giLandscapeChapterOpener replaces the heading with raw LaTeX)
% makes Quarto auto-insert an empty \chapter{} of its own — the book
% project model expects every chapter file to correspond to some
% heading, and falls back to a blank one when it finds none. Left
% alone, that produces exactly the wasted page this whole mechanism
% exists to avoid (an entirely blank "Chapter N" page, worse than the
% original problem since it has no title text at all), AND consumes
% a chapter number \giLandscapeChapterOpener's own \refstepcounter
% then increments past — which is why the very first test of this
% showed "Chapter 2" for what should have been Chapter 1. Fixed by
% making \chapter{} a complete no-op specifically when called with an
% empty title — real chapters (any actual title text) are completely
% unaffected; only Quarto's auto-inserted blank wrapper is silenced.
\makeatletter
\let\gi@origchapter\chapter
% Naively \renewcommand{\chapter}[1]{...} was the actual bug here —
% not just insufficient, but actively wrong. titlesec's real \chapter
% supports a starred form (\chapter*{...}, used everywhere in this
% project's own front/back matter — Copyright, Preface, Introduction,
% etc.) via \@ifstar-style star-detection built into how \chapter
% itself is defined. Overwriting \chapter with a plain, single-
% mandatory-argument macro destroys that: LaTeX's undelimited-
% argument grabbing takes whatever token comes right after the
% command name, so \chapter*{Copyright} — a single bare "*" character
% immediately following \chapter — had its "*" grabbed AS the
% argument, leaving "{Copyright}" to fall through as ordinary text,
% while the "*" (non-empty!) failed the empty-title check and got
% passed to \gi@origchapter as an UNSTARRED call — silently
% incrementing the chapter counter and mislabeling the title, for
% EVERY starred chapter in front and back matter, not just the one
% empty \chapter{} this was meant to catch. \@ifstar below restores
% correct dispatch: a starred call goes straight through to the real
% \chapter*, completely untouched by the empty-title logic, which
% only ever applies to the plain, unstarred form.
\renewcommand{\chapter}{\@ifstar{\gi@chapterstar}{\gi@chapternostar}}
\newcommand{\gi@chapterstar}[1]{\gi@origchapter*{#1}}
\newcommand{\gi@chapternostar}[1]{%
  \def\gi@chaptertitletest{#1}%
  \ifx\gi@chaptertitletest\@empty
  \else
    \gi@origchapter{#1}%
  \fi
}
\newcommand{\giLandscapeChapterOpener}[1]{%
  \clearpage
  \refstepcounter{chapter}
  \addcontentsline{toc}{chapter}{\protect\numberline{\thechapter}#1}%
  \markboth{\MakeUppercase{#1}}{\MakeUppercase{#1}}%
  \gdef\gi@landscapechaptertitle{#1}%
}
\newcommand{\giLandscapeChapterHeader}{%
  \thispagestyle{empty}%
  {\headingfont\Large\bfseries\color{giPrimary}%
   \chaptertitlename\ \thechapter\quad \gi@landscapechaptertitle%
   \hfill{\small\color{giSecondary}\giBookTitle}\par}%
  \vspace{4pt}{\color{giCrimson}\hrule height 0.6pt}\vspace{10pt}%
}
\makeatother

% ============================================================
% OPTIONAL LIST OF TABLES / LIST OF FIGURES
% A book with no captioned tables (or no captioned figures) shouldn't
% show an empty "List of Tables"/"List of Figures" page with a
% heading and nothing underneath it — that reads as broken, not as
% "there's nothing here." \@starttoc{lot}/\@starttoc{lof} alone can't
% do this: it doesn't know how many entries it's about to print until
% it's already printing them, and a Markdown heading (as used for the
% Table of Contents itself) forces its own chapter-opening page
% *unconditionally*, before we'd have any chance to check.
%
% \giOptionalTOCList{ext}{heading text} checks whether the relevant
% aux file (.lot or .lof) actually has at least one entry, and only
% prints the chapter heading + adds it to the Table of Contents if
% so. \@starttoc itself is always called unconditionally regardless —
% it's what (re-)opens that extension's write stream for the rest of
% the document, so table/figure captions later in the book have
% somewhere to record themselves. Making that call conditional too
% creates a permanent chicken-and-egg deadlock: a fresh build has no
% .lot/.lof file yet, so the check always finds nothing the first
% time, so \@starttoc never runs, so the write stream never opens, so
% nothing is ever recorded, on any pass, ever. When there genuinely
% are zero entries, \@starttoc silently produces no visible output on
% its own (nothing to read back), so calling it unconditionally is
% always safe.
%
% The "does it have entries" check reads the file back with
% \contentsline and \addvspace — the only two commands a stock
% .lot/.lof file contains — temporarily redefined to do nothing
% except flip a flag, so nothing in the file actually executes.
% \contentsline takes FOUR arguments here, not three: hyperref
% patches it to add a PDF anchor-name argument once loaded (which it
% always is here), and getting the arity wrong doesn't error — it
% silently leaves that fourth {argument} as leftover text, typeset
% literally on the page (hit exactly this: stray text like
% "figure.caption.9" printing where it shouldn't). The redefinitions
% are wrapped in their own local \makeatletter, separate from the one
% around this whole macro's definition — catcodes apply at the moment
% text is tokenized, not at macro-definition time, and the file's own
% content (e.g. \addvspace{10\p@}) is only tokenized when
% \InputIfFileExists actually reads it, much later, at this macro's
% call site, by which point the definition-time \makeatletter no
% longer applies unless it's re-established right here.
%
% (A per-caption counter, incremented from inside the \caption patch
% above instead of reading the .lot/.lof file, was tried first and
% abandoned: a raw-LaTeX table's \caption inside `longtable` doesn't
% reliably invoke the same \caption that patch redefines — `caption`
% package's `longtable` integration appears to bypass it — so table
% captions were never counted no matter how the counter logic was
% written. Reading the .lot/.lof file directly sidesteps this
% entirely: it checks the actual *result* of whatever mechanism wrote
% to that file, regardless of which internal LaTeX pathway did it.)
%
% Call it directly in raw LaTeX instead of a Markdown heading — see
% contents.qmd for the exact usage.
% ============================================================
\makeatletter
\newif\ifgi@tochasentries
\newcommand{\giOptionalTOCList}[2]{%
  \gi@tochasentriesfalse
  \begingroup
    \makeatletter
    \def\contentsline##1##2##3##4{\global\gi@tochasentriestrue}%
    \def\addvspace##1{}%
    \InputIfFileExists{\jobname.#1}{}{}%
  \endgroup
  \ifgi@tochasentries
    \chapter*{#2}
    \addcontentsline{toc}{chapter}{#2}
    \markboth{#2}{#2}
  \fi
  \@starttoc{#1}
}
\makeatother
